% Description of data/GZ/GZ_GOLD/GZ-GOLD_301.json
% Derived from Gajo & Barrón-Cedeño, "On Cross-Language Entity Label
% Projection and Recognition", CLiC-it 2024 (2024.clicit-1.47).
%
% This file is meant to be \input{} into a paper that already loads the
% booktabs package. It can also be compiled on its own by uncommenting the
% standalone preamble below.
%
% \documentclass{article}
% \usepackage{booktabs}
% \usepackage[T1]{fontenc}
% \begin{document}
% % Description of data/GZ/GZ_GOLD/GZ-GOLD_301.json
% Derived from Gajo & Barrón-Cedeño, "On Cross-Language Entity Label
% Projection and Recognition", CLiC-it 2024 (2024.clicit-1.47).
%
% This file is meant to be \input{} into a paper that already loads the
% booktabs package. It can also be compiled on its own by uncommenting the
% standalone preamble below.
%
% \documentclass{article}
% \usepackage{booktabs}
% \usepackage[T1]{fontenc}
% \begin{document}
% % Description of data/GZ/GZ_GOLD/GZ-GOLD_301.json
% Derived from Gajo & Barrón-Cedeño, "On Cross-Language Entity Label
% Projection and Recognition", CLiC-it 2024 (2024.clicit-1.47).
%
% This file is meant to be \input{} into a paper that already loads the
% booktabs package. It can also be compiled on its own by uncommenting the
% standalone preamble below.
%
% \documentclass{article}
% \usepackage{booktabs}
% \usepackage[T1]{fontenc}
% \begin{document}
% % Description of data/GZ/GZ_GOLD/GZ-GOLD_301.json
% Derived from Gajo & Barrón-Cedeño, "On Cross-Language Entity Label
% Projection and Recognition", CLiC-it 2024 (2024.clicit-1.47).
%
% This file is meant to be \input{} into a paper that already loads the
% booktabs package. It can also be compiled on its own by uncommenting the
% standalone preamble below.
%
% \documentclass{article}
% \usepackage{booktabs}
% \usepackage[T1]{fontenc}
% \begin{document}
% \input{dataset_description}
% \end{document}

\section{The GialloZafferano (GZ) Gold Dataset}
\label{sec:gz-gold}

The file \texttt{GZ-GOLD\_301.json} contains the gold portion of the
GialloZafferano (GZ) dataset, a manually annotated English--Italian parallel
corpus of culinary ingredient lists used to evaluate the cross-language entity
alignment and recognition models described in this work. The recipes were
collected from GialloZafferano,\footnote{\url{https://www.giallozafferano.it}}
where the English recipes are human translations of the Italian originals. The
data is partially out-of-domain with respect to the machine-translated training
data, as regards the food products, units of measure, and cooking processes it
covers, which makes it a suitable test bed for our models.

\paragraph{Format.}
The dataset is a JSON export from Label
Studio,\footnote{\url{https://labelstud.io}} the tool used for the manual
annotation. It is a list of $597$ recipe objects. Each object carries a unique
\texttt{id} together with a \texttt{data} field holding the parallel raw text
and metadata of the recipe, and an \texttt{annotations} field holding the
manual labels. The \texttt{data} field provides, for each of the two languages
$\ell \in \{\texttt{en}, \texttt{it}\}$, the source URL (\texttt{url\_\ell}),
the recipe title (\texttt{title\_\ell}), a short presentation
(\texttt{presentation\_\ell}), the ingredient list (\texttt{ingredients\_\ell}),
and the preparation steps (\texttt{preparation\_\ell}). The annotations are
applied to the \texttt{ingredients\_en} and \texttt{ingredients\_it} strings,
in which the individual ingredients are separated by semicolons.

\paragraph{Annotations.}
The corpus was annotated by a single professional translator, a native speaker
of Italian holding an MA in Specialized Translation in English and Spanish, in
two successive tasks. The first is a multi-class span-annotation (NER) task in
which spans of the ingredient lists, in both languages, are assigned one of the
recipe Named Entity (r-NE) categories: \textsc{food}, \textsc{quantity},
\textsc{unit}, \textsc{process}, \textsc{physical\_quality}, \textsc{color},
\textsc{taste}, \textsc{purpose}, and \textsc{part}.\footnote{A handful of spans
($8$ in English, $33$ in Italian) carry the residual \textsc{quality/brand}
label, which is not part of the main label scheme.} The second is a
cross-language entity-linking task in which each source (English) entity is
linked to its corresponding target (Italian) entity, with the possibility of a
source entity being linked to more than one target entity and vice versa. The
latter situation arises because some recipes list more than one ingredient
option in one language but not in the other, e.g.\ \emph{Cocomero (anguria) 1
fetta} vs.\ \emph{Watermelon 1 slice}.

In the JSON, each annotated span is a \texttt{result} element of
\texttt{type}~\texttt{labels}, with \texttt{from\_name} either \texttt{label\_en}
or \texttt{label\_it}, a \texttt{value} giving the character-level
\texttt{start}/\texttt{end} offsets, the surface \texttt{text}, and the assigned
\texttt{labels}, plus a unique span \texttt{id}. Each alignment is a
\texttt{result} element of \texttt{type}~\texttt{relation} connecting a source
and a target span through their \texttt{id}s (\texttt{from\_id}/\texttt{to\_id}),
with a bidirectional \texttt{direction}.

\paragraph{Coverage.}
All $597$ recipes are bilingual and carry NER annotations in both languages.
The cross-language alignments were annotated on a subset of $300$ recipes.
Because the paired recipes do not always coincide completely---some ingredients
may be missing in one language or replaced by an equivalent rather than the same
food product---alignments were not annotated whenever the share of source
ingredients missing from the target recipe exceeded a heuristic threshold of
$1/3$, so as to avoid aligning excessively divergent recipes. In total, the
dataset contains $46{,}903$ entity annotations ($26{,}631$ in English and
$20{,}272$ in Italian) and $9{,}842$ alignments (on average $\approx 33$ per
aligned recipe). Quantities and units of measure are localized: the Italian
recipes use SI units whereas the English ones additionally report imperial
units, which is reflected in the considerably lower number of
\textsc{quantity} and \textsc{unit} spans on the Italian side.
Table~\ref{tab:gz-classes} reports the per-class distribution of the entity
annotations.

\begin{table}[t]
  \centering
  \begin{tabular}{lrr}
    \toprule
    Class & GZ (en) & GZ (it) \\
    \midrule
    food         & 5{,}958  & 6{,}473  \\
    quantity     & 10{,}186 & 6{,}564  \\
    unit         & 8{,}148  & 4{,}450  \\
    process      & 217      & 265      \\
    physical q.  & 1{,}245  & 1{,}547  \\
    color        & 482      & 479      \\
    taste        & 98       & 72       \\
    purpose      & 69       & 126      \\
    part         & 220      & 263      \\
    \midrule
    total        & 26{,}631 & 20{,}272 \\
    \bottomrule
  \end{tabular}
  \caption{Per-class distribution of the entity annotations in the GZ gold
  dataset (\texttt{GZ-GOLD\_301.json}), for the English and Italian portions.}
  \label{tab:gz-classes}
\end{table}

% \end{document}

\section{The GialloZafferano (GZ) Gold Dataset}
\label{sec:gz-gold}

The file \texttt{GZ-GOLD\_301.json} contains the gold portion of the
GialloZafferano (GZ) dataset, a manually annotated English--Italian parallel
corpus of culinary ingredient lists used to evaluate the cross-language entity
alignment and recognition models described in this work. The recipes were
collected from GialloZafferano,\footnote{\url{https://www.giallozafferano.it}}
where the English recipes are human translations of the Italian originals. The
data is partially out-of-domain with respect to the machine-translated training
data, as regards the food products, units of measure, and cooking processes it
covers, which makes it a suitable test bed for our models.

\paragraph{Format.}
The dataset is a JSON export from Label
Studio,\footnote{\url{https://labelstud.io}} the tool used for the manual
annotation. It is a list of $597$ recipe objects. Each object carries a unique
\texttt{id} together with a \texttt{data} field holding the parallel raw text
and metadata of the recipe, and an \texttt{annotations} field holding the
manual labels. The \texttt{data} field provides, for each of the two languages
$\ell \in \{\texttt{en}, \texttt{it}\}$, the source URL (\texttt{url\_\ell}),
the recipe title (\texttt{title\_\ell}), a short presentation
(\texttt{presentation\_\ell}), the ingredient list (\texttt{ingredients\_\ell}),
and the preparation steps (\texttt{preparation\_\ell}). The annotations are
applied to the \texttt{ingredients\_en} and \texttt{ingredients\_it} strings,
in which the individual ingredients are separated by semicolons.

\paragraph{Annotations.}
The corpus was annotated by a single professional translator, a native speaker
of Italian holding an MA in Specialized Translation in English and Spanish, in
two successive tasks. The first is a multi-class span-annotation (NER) task in
which spans of the ingredient lists, in both languages, are assigned one of the
recipe Named Entity (r-NE) categories: \textsc{food}, \textsc{quantity},
\textsc{unit}, \textsc{process}, \textsc{physical\_quality}, \textsc{color},
\textsc{taste}, \textsc{purpose}, and \textsc{part}.\footnote{A handful of spans
($8$ in English, $33$ in Italian) carry the residual \textsc{quality/brand}
label, which is not part of the main label scheme.} The second is a
cross-language entity-linking task in which each source (English) entity is
linked to its corresponding target (Italian) entity, with the possibility of a
source entity being linked to more than one target entity and vice versa. The
latter situation arises because some recipes list more than one ingredient
option in one language but not in the other, e.g.\ \emph{Cocomero (anguria) 1
fetta} vs.\ \emph{Watermelon 1 slice}.

In the JSON, each annotated span is a \texttt{result} element of
\texttt{type}~\texttt{labels}, with \texttt{from\_name} either \texttt{label\_en}
or \texttt{label\_it}, a \texttt{value} giving the character-level
\texttt{start}/\texttt{end} offsets, the surface \texttt{text}, and the assigned
\texttt{labels}, plus a unique span \texttt{id}. Each alignment is a
\texttt{result} element of \texttt{type}~\texttt{relation} connecting a source
and a target span through their \texttt{id}s (\texttt{from\_id}/\texttt{to\_id}),
with a bidirectional \texttt{direction}.

\paragraph{Coverage.}
All $597$ recipes are bilingual and carry NER annotations in both languages.
The cross-language alignments were annotated on a subset of $300$ recipes.
Because the paired recipes do not always coincide completely---some ingredients
may be missing in one language or replaced by an equivalent rather than the same
food product---alignments were not annotated whenever the share of source
ingredients missing from the target recipe exceeded a heuristic threshold of
$1/3$, so as to avoid aligning excessively divergent recipes. In total, the
dataset contains $46{,}903$ entity annotations ($26{,}631$ in English and
$20{,}272$ in Italian) and $9{,}842$ alignments (on average $\approx 33$ per
aligned recipe). Quantities and units of measure are localized: the Italian
recipes use SI units whereas the English ones additionally report imperial
units, which is reflected in the considerably lower number of
\textsc{quantity} and \textsc{unit} spans on the Italian side.
Table~\ref{tab:gz-classes} reports the per-class distribution of the entity
annotations.

\begin{table}[t]
  \centering
  \begin{tabular}{lrr}
    \toprule
    Class & GZ (en) & GZ (it) \\
    \midrule
    food         & 5{,}958  & 6{,}473  \\
    quantity     & 10{,}186 & 6{,}564  \\
    unit         & 8{,}148  & 4{,}450  \\
    process      & 217      & 265      \\
    physical q.  & 1{,}245  & 1{,}547  \\
    color        & 482      & 479      \\
    taste        & 98       & 72       \\
    purpose      & 69       & 126      \\
    part         & 220      & 263      \\
    \midrule
    total        & 26{,}631 & 20{,}272 \\
    \bottomrule
  \end{tabular}
  \caption{Per-class distribution of the entity annotations in the GZ gold
  dataset (\texttt{GZ-GOLD\_301.json}), for the English and Italian portions.}
  \label{tab:gz-classes}
\end{table}

% \end{document}

\section{The GialloZafferano (GZ) Gold Dataset}
\label{sec:gz-gold}

The file \texttt{GZ-GOLD\_301.json} contains the gold portion of the
GialloZafferano (GZ) dataset, a manually annotated English--Italian parallel
corpus of culinary ingredient lists used to evaluate the cross-language entity
alignment and recognition models described in this work. The recipes were
collected from GialloZafferano,\footnote{\url{https://www.giallozafferano.it}}
where the English recipes are human translations of the Italian originals. The
data is partially out-of-domain with respect to the machine-translated training
data, as regards the food products, units of measure, and cooking processes it
covers, which makes it a suitable test bed for our models.

\paragraph{Format.}
The dataset is a JSON export from Label
Studio,\footnote{\url{https://labelstud.io}} the tool used for the manual
annotation. It is a list of $597$ recipe objects. Each object carries a unique
\texttt{id} together with a \texttt{data} field holding the parallel raw text
and metadata of the recipe, and an \texttt{annotations} field holding the
manual labels. The \texttt{data} field provides, for each of the two languages
$\ell \in \{\texttt{en}, \texttt{it}\}$, the source URL (\texttt{url\_\ell}),
the recipe title (\texttt{title\_\ell}), a short presentation
(\texttt{presentation\_\ell}), the ingredient list (\texttt{ingredients\_\ell}),
and the preparation steps (\texttt{preparation\_\ell}). The annotations are
applied to the \texttt{ingredients\_en} and \texttt{ingredients\_it} strings,
in which the individual ingredients are separated by semicolons.

\paragraph{Annotations.}
The corpus was annotated by a single professional translator, a native speaker
of Italian holding an MA in Specialized Translation in English and Spanish, in
two successive tasks. The first is a multi-class span-annotation (NER) task in
which spans of the ingredient lists, in both languages, are assigned one of the
recipe Named Entity (r-NE) categories: \textsc{food}, \textsc{quantity},
\textsc{unit}, \textsc{process}, \textsc{physical\_quality}, \textsc{color},
\textsc{taste}, \textsc{purpose}, and \textsc{part}.\footnote{A handful of spans
($8$ in English, $33$ in Italian) carry the residual \textsc{quality/brand}
label, which is not part of the main label scheme.} The second is a
cross-language entity-linking task in which each source (English) entity is
linked to its corresponding target (Italian) entity, with the possibility of a
source entity being linked to more than one target entity and vice versa. The
latter situation arises because some recipes list more than one ingredient
option in one language but not in the other, e.g.\ \emph{Cocomero (anguria) 1
fetta} vs.\ \emph{Watermelon 1 slice}.

In the JSON, each annotated span is a \texttt{result} element of
\texttt{type}~\texttt{labels}, with \texttt{from\_name} either \texttt{label\_en}
or \texttt{label\_it}, a \texttt{value} giving the character-level
\texttt{start}/\texttt{end} offsets, the surface \texttt{text}, and the assigned
\texttt{labels}, plus a unique span \texttt{id}. Each alignment is a
\texttt{result} element of \texttt{type}~\texttt{relation} connecting a source
and a target span through their \texttt{id}s (\texttt{from\_id}/\texttt{to\_id}),
with a bidirectional \texttt{direction}.

\paragraph{Coverage.}
All $597$ recipes are bilingual and carry NER annotations in both languages.
The cross-language alignments were annotated on a subset of $300$ recipes.
Because the paired recipes do not always coincide completely---some ingredients
may be missing in one language or replaced by an equivalent rather than the same
food product---alignments were not annotated whenever the share of source
ingredients missing from the target recipe exceeded a heuristic threshold of
$1/3$, so as to avoid aligning excessively divergent recipes. In total, the
dataset contains $46{,}903$ entity annotations ($26{,}631$ in English and
$20{,}272$ in Italian) and $9{,}842$ alignments (on average $\approx 33$ per
aligned recipe). Quantities and units of measure are localized: the Italian
recipes use SI units whereas the English ones additionally report imperial
units, which is reflected in the considerably lower number of
\textsc{quantity} and \textsc{unit} spans on the Italian side.
Table~\ref{tab:gz-classes} reports the per-class distribution of the entity
annotations.

\begin{table}[t]
  \centering
  \begin{tabular}{lrr}
    \toprule
    Class & GZ (en) & GZ (it) \\
    \midrule
    food         & 5{,}958  & 6{,}473  \\
    quantity     & 10{,}186 & 6{,}564  \\
    unit         & 8{,}148  & 4{,}450  \\
    process      & 217      & 265      \\
    physical q.  & 1{,}245  & 1{,}547  \\
    color        & 482      & 479      \\
    taste        & 98       & 72       \\
    purpose      & 69       & 126      \\
    part         & 220      & 263      \\
    \midrule
    total        & 26{,}631 & 20{,}272 \\
    \bottomrule
  \end{tabular}
  \caption{Per-class distribution of the entity annotations in the GZ gold
  dataset (\texttt{GZ-GOLD\_301.json}), for the English and Italian portions.}
  \label{tab:gz-classes}
\end{table}

% \end{document}

\section{The GialloZafferano (GZ) Gold Dataset}
\label{sec:gz-gold}

The file \texttt{GZ-GOLD\_301.json} contains the gold portion of the
GialloZafferano (GZ) dataset, a manually annotated English--Italian parallel
corpus of culinary ingredient lists used to evaluate the cross-language entity
alignment and recognition models described in this work. The recipes were
collected from GialloZafferano,\footnote{\url{https://www.giallozafferano.it}}
where the English recipes are human translations of the Italian originals. The
data is partially out-of-domain with respect to the machine-translated training
data, as regards the food products, units of measure, and cooking processes it
covers, which makes it a suitable test bed for our models.

\paragraph{Format.}
The dataset is a JSON export from Label
Studio,\footnote{\url{https://labelstud.io}} the tool used for the manual
annotation. It is a list of $597$ recipe objects. Each object carries a unique
\texttt{id} together with a \texttt{data} field holding the parallel raw text
and metadata of the recipe, and an \texttt{annotations} field holding the
manual labels. The \texttt{data} field provides, for each of the two languages
$\ell \in \{\texttt{en}, \texttt{it}\}$, the source URL (\texttt{url\_\ell}),
the recipe title (\texttt{title\_\ell}), a short presentation
(\texttt{presentation\_\ell}), the ingredient list (\texttt{ingredients\_\ell}),
and the preparation steps (\texttt{preparation\_\ell}). The annotations are
applied to the \texttt{ingredients\_en} and \texttt{ingredients\_it} strings,
in which the individual ingredients are separated by semicolons.

\paragraph{Annotations.}
The corpus was annotated by a single professional translator, a native speaker
of Italian holding an MA in Specialized Translation in English and Spanish, in
two successive tasks. The first is a multi-class span-annotation (NER) task in
which spans of the ingredient lists, in both languages, are assigned one of the
recipe Named Entity (r-NE) categories: \textsc{food}, \textsc{quantity},
\textsc{unit}, \textsc{process}, \textsc{physical\_quality}, \textsc{color},
\textsc{taste}, \textsc{purpose}, and \textsc{part}.\footnote{A handful of spans
($8$ in English, $33$ in Italian) carry the residual \textsc{quality/brand}
label, which is not part of the main label scheme.} The second is a
cross-language entity-linking task in which each source (English) entity is
linked to its corresponding target (Italian) entity, with the possibility of a
source entity being linked to more than one target entity and vice versa. The
latter situation arises because some recipes list more than one ingredient
option in one language but not in the other, e.g.\ \emph{Cocomero (anguria) 1
fetta} vs.\ \emph{Watermelon 1 slice}.

In the JSON, each annotated span is a \texttt{result} element of
\texttt{type}~\texttt{labels}, with \texttt{from\_name} either \texttt{label\_en}
or \texttt{label\_it}, a \texttt{value} giving the character-level
\texttt{start}/\texttt{end} offsets, the surface \texttt{text}, and the assigned
\texttt{labels}, plus a unique span \texttt{id}. Each alignment is a
\texttt{result} element of \texttt{type}~\texttt{relation} connecting a source
and a target span through their \texttt{id}s (\texttt{from\_id}/\texttt{to\_id}),
with a bidirectional \texttt{direction}.

\paragraph{Coverage.}
All $597$ recipes are bilingual and carry NER annotations in both languages.
The cross-language alignments were annotated on a subset of $300$ recipes.
Because the paired recipes do not always coincide completely---some ingredients
may be missing in one language or replaced by an equivalent rather than the same
food product---alignments were not annotated whenever the share of source
ingredients missing from the target recipe exceeded a heuristic threshold of
$1/3$, so as to avoid aligning excessively divergent recipes. In total, the
dataset contains $46{,}903$ entity annotations ($26{,}631$ in English and
$20{,}272$ in Italian) and $9{,}842$ alignments (on average $\approx 33$ per
aligned recipe). Quantities and units of measure are localized: the Italian
recipes use SI units whereas the English ones additionally report imperial
units, which is reflected in the considerably lower number of
\textsc{quantity} and \textsc{unit} spans on the Italian side.
Table~\ref{tab:gz-classes} reports the per-class distribution of the entity
annotations.

\begin{table}[t]
  \centering
  \begin{tabular}{lrr}
    \toprule
    Class & GZ (en) & GZ (it) \\
    \midrule
    food         & 5{,}958  & 6{,}473  \\
    quantity     & 10{,}186 & 6{,}564  \\
    unit         & 8{,}148  & 4{,}450  \\
    process      & 217      & 265      \\
    physical q.  & 1{,}245  & 1{,}547  \\
    color        & 482      & 479      \\
    taste        & 98       & 72       \\
    purpose      & 69       & 126      \\
    part         & 220      & 263      \\
    \midrule
    total        & 26{,}631 & 20{,}272 \\
    \bottomrule
  \end{tabular}
  \caption{Per-class distribution of the entity annotations in the GZ gold
  dataset (\texttt{GZ-GOLD\_301.json}), for the English and Italian portions.}
  \label{tab:gz-classes}
\end{table}
